% ===== 模型检验 =====
\section{模型检验}

\subsection{误差分析}
% 数值验证：模型结果与实际数据/标准解的对比

% 对比表使用 booktabs 三线表：
% \begin{table}[!htbp]
%   \centering
%   \caption{模型误差分析结果}
%   \label{tab:error}
%   \begin{tabular}{cccc}
%     \toprule[1.5pt]
%     指标 & 训练集 & 测试集 & 说明 \\
%     \midrule[1pt]
%     $R^2$ & ... & ... & 决定系数 \\
%     RMSE & ... & ... & 均方根误差 \\
%     MAE & ... & ... & 平均绝对误差 \\
%     MAPE(\%) & ... & ... & 平均绝对百分比误差 \\
%     \bottomrule[1.5pt]
%   \end{tabular}
% \end{table}

\subsection{敏感性分析}
% 参考 references/敏感性分析写作模板.md 的标准结构
% 每个主要模型必须有敏感性分析章节
% 结构：参数选取 → 分析方法（OAT ±10%/±20%/±30%）→ 结果与分析（龙卷风图/敏感性曲线）→ 稳健性结论
% 结论须量化："参数X变化±20%，目标值变化±Y%"

% 龙卷风图示例标签：
% \label{fig:sensitivity}
